\section{Introduction}

\label{sec:introduction}
% ══════════════════════════════════════════════════════════════════════════════

Citizen science---the participation of non-professional scientists in
research data collection---has transformed biodiversity monitoring.
eBird receives over 200 million bird observations per year
\citep{sullivan2014ebird}. iNaturalist has documented over 150 million
observations across all taxa \citep{inaturalist2024}. These platforms
generate data at scales that professional surveys cannot match, and
citizen science observations now underpin a substantial fraction of
published biodiversity research \citep{theobald2015global}.

Despite these achievements, citizen science data has well-documented
quality and bias issues:

\begin{enumerate}[label=(\roman*)]
  \item \textbf{Spatial bias}: Observations cluster near roads,
    population centers, and national parks, leaving vast areas
    unsampled \citep{boakes2010distorted}.
  \item \textbf{Temporal bias}: Weekend and holiday peaks produce
    uneven temporal coverage, confounding phenological analyses
    \citep{courter2013weekend}.
  \item \textbf{Taxonomic bias}: Charismatic megafauna and birds are
    overrepresented; invertebrates, fungi, and cryptic species are
    underrepresented \citep{troudet2017taxonomic}.
  \item \textbf{Observer heterogeneity}: Skill levels range from novice
    to expert, producing variable identification accuracy
    \citep{kelling2015can}.
  \item \textbf{Motivation decay}: Most contributors submit a few
    observations then disengage, creating a heavy-tailed contribution
    distribution \citep{sauermann2015crowd}.
\end{enumerate}

\subsection{The Gamification Opportunity}

Gamification---applying game design elements to non-game contexts---has
shown promise in sustaining engagement and directing effort in citizen
science \citep{bowser2013gamifying, eveleigh2014designing}. However,
naive gamification (leaderboards, badges) can exacerbate biases:
competitive leaderboards incentivize easy observations near home, not
scientifically valuable observations in undersampled areas
\citep{crowston2020citizen}.

\subsection{Contributions}

We present ZooPlatform with the following contributions:

\begin{enumerate}
  \item An \textbf{adaptive gamification engine} powered by
    reinforcement learning that personalizes challenges to direct
    observer effort toward scientifically valuable observations,
    reducing sampling bias by 47\%.

  \item An \textbf{AI-assisted identification pipeline} integrating
    TaxoNet (ZOO-CLS-2026) for real-time species suggestions with
    confidence-calibrated outputs.

  \item A \textbf{Bayesian data quality model} that estimates
    per-observation reliability by jointly modeling observer expertise,
    AI confidence, environmental context, and community validation.

  \item A \textbf{token-based incentive system} where verified
    contributions earn ZooRep tokens that integrate with Zoo DAO
    governance, creating a direct link between scientific contribution
    and decision-making power.

  \item \textbf{Empirical evaluation} from 14 months of platform
    operation with 23{,}847 contributors and 4.2 million observations.
\end{enumerate}


% ══════════════════════════════════════════════════════════════════════════════
